\section{Separate inner precision and calculus of flat sectors}
\label{sec:flat-sector-operations}

A finite calculus for the scale of Section~\ref{sec:flat-sectors}
uses a common positive monomial target coordinate
\[
 u=\left(\frac{y-y_0}{a}\right)^{1/q},\qquad
 E=e^{-c/u^p},\qquad p,c>0,
\]
for every operand. The phase and chart are part of the scale; combining
different phases requires an additional change-of-scale argument.

The asymptotic representation distinguishes two errors:
\begin{equation}
 S(y)=C_0(u)+\sum_{k=1}^{N}E^k
 \left(T_k(u)+\OO(u^{\rho_k}\M(u)^{d_k})\right)
 +\OO(E^{N+1}u^{\rho_*}\M(u)^{d_*}).
 \label{eq:flat-two-precision}
\end{equation}
Here $C_0$ is exact, each $T_k$ is a finite power--log expression, and
$\rho_k=+\infty$ means that the coefficient is exact. The remainder
is a sum with one term for each nonzero inner error and a separate term
for the complete omitted exponential tail. These are asymptotic bounds for
fixed data, whose constants and threshold are existential.

\subsection{Inner truncation and multiplication}

An inner truncation at $h$ retains powers strictly below $h$ in every
positive sector. The zero sector remains exact, including any of its powers
at or above $h$. If a coefficient loses its first term at power $\alpha$,
its inner error has power $\alpha$ and the degree of the corresponding
logarithmic polynomial. An existing, less precise bound is preserved.
Increasing the formal cutoff after truncation does not recover discarded
coefficient information.

For example, truncating the inner coefficients of
\[
 S=y-y^2E+(y^2+2y^3)E^2+\OO(y^2E^3),\qquad E=e^{-1/y},
\]
at $h=3$ gives
\[
 S=y-y^2E+y^2E^2+E^2\OO(y^3)+\OO(y^2E^3).
\]
The error $E^2\OO(y^3)$ cannot be put into the third sector: its ratio to
any $E^3$ times a fixed power--log expression is unbounded as $y\to0^+$.

Multiplication convolves exponential degrees and uses the ordinary
precision-tracked power--log multiplication inside each degree. If
$A=\sum_{i\le N_A}E^i A_i+R_A$ and
$B=\sum_{j\le N_B}E^j B_j+R_B$, the result retains degrees through
$N=\min(N_A,N_B)$. Its tail includes all three kinds of contributions:
\[
 \sum_{i+j>N}E^{i+j}A_iB_j,
 \qquad R_A\sum_j E^j B_j+R_B\sum_i E^i A_i,
 \qquad R_A R_B.
\]
Each coefficient bound includes its own inner uncertainty. Since $0<E<1$,
a term in degree $k>N$ can conservatively be bounded in degree $N+1$
while retaining its power--log envelope. This can give a looser algebraic
power in the final sector bound; it is always a valid bound. In particular,
unknown tails are added by bounds and never cancelled because some displayed
coefficients cancel.

Polynomial composition $H(S)$ follows from addition and multiplication.
The coefficients of $H$ may be exact finite real power--log expressions
in the target chart. Horner's formula computes the composition through
successive products and gives the exact zero sector $H(C_0)$. A constant
observable has no dependence on the unknown input tail; the zero observable
therefore has zero remainder. General nonpolynomial composition and changes
of phase require additional analytic and scale-compatibility arguments.

\subsection{Qualified differentiation}

A magnitude bound alone does not justify its differentiated bound.
For the exact flat inverse, the uniform holomorphic implicit-function
construction supplies bounds for every fixed derivative order. The exact
finite forward model is essential to this argument. Addition,
multiplication, polynomial composition, and inner truncation preserve the
required local holomorphic estimates. A general real-axis remainder has
no such derivative bound without a separate regularity hypothesis.

For a retained monomial the target derivative is explicit:
\begin{equation}
\begin{aligned}
 \frac{d}{dy}\left(E^k u^\alpha P(\log u)\right)
 =\frac{E^k}{a q}\Bigl[&u^{\alpha-q}
       \bigl(\alpha P(\log u)+P'(\log u)\bigr)\\
       &+kcp\,u^{\alpha-p-q}P(\log u)\Bigr].
\end{aligned}
 \label{eq:flat-target-derivative}
\end{equation}
Thus the exponential degree stays $k$. For $k>0$, a qualified inner
bound at $(\rho_k,d_k)$ is transported to
$(\rho_k-p-q,d_k)$. When known differentiated terms are then discarded at
that same boundary power, their logarithmic degree can increase the final
inner degree bound. The complete tail has the same power shift
$\rho_*\mapsto\rho_*-p-q$. This formula applies also when $q<0$;
the shift is a signed exact number, not its absolute value.

For completeness, the analytic tail estimate supports this derivative
claim. Fix a small positive real $u$ and use a complex disk
\[
 |v-u|<\eta u^{p+1},
\]
with fixed sufficiently small $\eta>0$. On this disk the chosen local
power and logarithm branches are holomorphic, their power--log envelopes
are comparable to those at $u$, and
\[
 v^{-p}-u^{-p}=\OO(1).
\]
Consequently $e^{-c/v^p}$ differs from $e^{-c/u^p}$ by a bounded factor.
The normalized implicit-function construction in
Section~\ref{sec:flat-sectors} extends uniformly to these disks. Its
geometric tail estimate therefore holds there with changed constants.
Cauchy's estimate costs a factor $u^{-p-1}$ for one source derivative.
The chain-rule factor $du/dy=(a q)^{-1}u^{1-q}$ gives the asserted loss
$u^{-p-q}$ for one target derivative. Every fixed higher derivative is
handled on nested disks, or by the corresponding target disk of radius
comparable to $u^{p+q}$; constants may depend on the derivative order.
No differentiation of an arbitrary real-axis Big-O bound is used.

The finitely many discarded inner terms are themselves analytic
power--log functions, so the same bound, and the exact phase derivative,
apply to their errors. After differentiation, terms are retained only within the transported precision. For the previous example, the derivative
of the truncated representation has the valid form
\[
 S'=1-(1+2y)E+2E^2+E^2\OO(y)+\OO(E^3),
\]
where the last bound is the transported $\OO(y^2E^3)$.
The full second-sector derivative contains further powers of $y$ that the
truncated data do not determine.

The two levels of precision remain distinct under all these operations.
A power-sized uncertainty in sector $k$ cannot be replaced by a tail in
sector $k+1$, even if many known coefficient terms cancel. The analytic
regularity required for differentiation is likewise separate from the
magnitude of the remainder. These distinctions prevent formal coefficient
arithmetic from asserting precision not present in the original data.
